\chapter{Le premier palier}
\sessioninfo{Session 11}{Mardi 17 novembre 2026}

Au pied de Firefinger, la première échelle disparaissait dans la nuit. Dix-huit mètres de barreaux séparaient le groupe d'une corniche à peine visible, tandis que la flamme du sommet et la lune dessinaient de faibles reflets sur la roche. Morrÿs observa une dernière fois ses compagnons avant l'ascension. Azaka transpirait abondamment et respirait bien plus vite qu'à l'ordinaire.

La guide finit par reconnaître qu'elle supportait mal la hauteur. Elle aurait volontiers défendu les qualités du sol ferme, mais son masque l'attendait quelque part au-dessus d'eux. Firefinger s'élevait sur près de quatre-vingt-dix mètres et ce premier obstacle ne représentait qu'une fraction de la montée. Il faudrait pourtant commencer par lui.

Destos suggéra de diviser le groupe, les plus discrets partant en reconnaissance pendant que les autres resteraient en contrebas. L'idée fut rapidement abandonnée : ceux qui atteindraient la corniche se retrouveraient seuls en cas d'attaque, tandis que les autres risqueraient d'attirer les créatures en restant au pied de la tour. Ils monteraient donc un par un, chacun couvrant le suivant. Kael passerait en tête, suivi de Morrÿs, d'Azaka, d'Obélik et enfin de Destos.

Avant de partir, Kael convint de deux signes avec ses compagnons. Un cercle formé avec les mains signifierait que la voie était libre ; une large croix leur ordonnerait de rester en bas. Il saisit les premiers barreaux et s'éleva sans bruit. À mesure qu'il montait, le vent se renforçait dans son dos et le sol s'éloignait sous ses pieds.

La corniche donnait sur l'entrée d'une grotte où la lumière extérieure ne pénétrait presque pas. Grâce à sa vision dans le noir, Kael distingua un mur à plusieurs mètres et, juste avant, une rupture dans le sol. Il tendit l'oreille, ne perçut aucun mouvement, puis forma le cercle convenu. Morrÿs pouvait monter.

Le nain le rejoignit et posa ses mains contre la pierre. Pendant quelques instants, il étendit ses sens dans la roche à la recherche de vibrations ou d'un déplacement autour d'eux. Rien ne lui répondit. Azaka commença alors son ascension, tremblant à chaque barreau, avant de se hisser sur la corniche et de s'allonger aussitôt contre le sol.

Obélik arriva ensuite sans avoir allumé la torche rangée dans son sac. Il ne voyait rien dans la grotte et préféra attendre plutôt que de révéler leur présence à toute la tour. Destos, dernier à gravir l'échelle, fit briller son symbole sacré sous ses vêtements afin de ne laisser filtrer qu'une faible lueur. Lorsqu'il rejoignit le goliath, celui-ci retrouva enfin l'usage de ses yeux.

\begin{narrateur}
Destos venait de trouver sa véritable vocation : torche humaine, avec variateur intégré.
\end{narrateur}

\scenebreak

La lumière révéla une fosse profonde d'environ trois mètres au milieu de la cavité. Une partie du groupe descendit pour examiner les lieux, tandis que Kael demeurait sur le rebord, son arc à portée de main. À peine leurs pieds touchèrent-ils le fond que Morrÿs perçut plusieurs vibrations dans la roche. Quatre araignées jaillirent aussitôt des anfractuosités et se ruèrent sur les intrus.

La première morsure manqua Destos, mais une seconde créature bondit à sa gorge et planta ses crochets dans son cou. Morrÿs fut touché à son tour ; sa résistance naturelle atténua le poison sans lui épargner la douleur. Trois araignées convergèrent presque aussitôt vers Obélik. Le goliath durcit sa peau au moment de l'impact et repoussa une première attaque, mais les suivantes traversèrent sa défense.

\illustration[0.95\textwidth]{2026-11-17/embuscade-araignees.png}
{\textit{Dans la fosse du premier palier, quatre araignées surgirent autour des aventuriers.}}
{embuscade_araignees_firefinger}

Destos réagit en faisant jaillir la lumière qui l'habitait. Il rendit une partie de ses forces à Obélik, puis referma sa propre blessure d'une courte prière. Le goliath, lui, cessa de réfléchir. Sa colère éclata, sa hache fendit l'air et la première araignée mourut sous le choc.

Azaka frappa une autre créature à l'épée pendant que Kael, resté en hauteur, déplaçait sa marque sur sa nouvelle proie. Il banda son arc et abattit une araignée d'une flèche. Morrÿs tenta de lancer un \textit{Fire Bolt}, mais la proximité des monstres gêna son geste et le projectile se perdit contre la roche. Il recula derrière Obélik pour retrouver de l'espace.

Les deux araignées à sa portée profitèrent de son mouvement. Leurs crochets l'atteignirent avant qu'il ne puisse se mettre à l'abri, ne lui laissant plus que quelques forces. Destos fit résonner un glas funèbre dans la fosse ; la chair d'une des créatures commença à se flétrir sous l'effet de sa magie. Obélik leva alors sa grande hache pour l'achever.

La lame heurta la roche et lui échappa.

Morrÿs, placé derrière lui, reçut la hache de plein fouet et s'effondra.

Pendant que le nain gisait inconscient, Azaka décocha une flèche qui abattit une troisième araignée. Elle sauta ensuite dans la fosse, tira son épée et blessa la dernière. Kael tenta de l'aider depuis le rebord, mais son tir passa à côté. La créature se jeta alors sur Azaka et planta ses crochets dans son bras.

Tous virent le poison pénétrer dans la plaie. Pourtant, lorsque l'araignée retira ses crochets, la blessure se referma aussitôt sans laisser couler une seule goutte de sang. Kael reconnut le phénomène observé après leur premier affrontement sur la rivière. Cette fois, personne ne pouvait prétendre ne pas l'avoir vu.

Destos se précipita auprès de Morrÿs et stabilisa son état avant que la mort ne puisse l'emporter. Obélik, toujours animé par sa rage, ne prit même pas le temps de ramasser sa hache. Il bondit sur la dernière araignée et lui asséna un coup de poing si violent que sa mâchoire s'enfonça dans son crâne. Le silence retomba enfin sur la fosse.

\scenebreak

Morrÿs reprit conscience quelques minutes plus tard, encore sonné par ce qu'il venait de subir. Il but la potion de soin que Yarno lui avait remise avant leur départ, tandis qu'Obélik utilisait la sienne pour apaiser les nombreuses morsures reçues pendant l'affrontement. Azaka leur assura qu'elle disposait également de quelques réserves, mais tous préférèrent les conserver pour la suite. Ils ignoraient combien de temps les séparait encore d'un véritable repos.

Kael récupéra deux de ses flèches parmi les cadavres. La fosse ne contenait rien d'autre d'utile, pas davantage que les crevasses alentour. Sa plaie refermée, Azaka ne donna aucune nouvelle explication sur ce qui venait de se produire. Le groupe remit donc cette question à plus tard et chercha comment poursuivre l'ascension.

Des lianes et de nombreuses prises permettaient d'atteindre sans grande difficulté le rebord opposé, trois mètres plus haut. Azaka le franchit presque d'un bond, avec une agilité féline, puis les autres la rejoignirent. Destos demeura un moment en contrebas pour parcourir la fosse et rechercher la moindre trace de magie. Son rituel ne révéla aucun passage caché ni aucun objet enchanté.

La petite salle supérieure n'offrait que de la roche, des crevasses et une corde suspendue en son centre. Celle-ci disparaissait dans un conduit vertical haut d'environ vingt-cinq mètres, bien au-delà de la première échelle. Azaka leur expliqua que Firefinger avait autrefois servi de tour de guet aux habitants de Chult. Un unique chemin permettait de contrôler les arrivées et, si nécessaire, de refouler facilement ceux qui tentaient de monter.

Obélik éprouva la corde de tout son poids. Elle semblait suffisamment solide, ce qui ne la rendait pas plus accueillante. Kael tenta le premier de s'y hisser, mais ses mains glissèrent presque immédiatement. Avec l'aide du goliath, il gagna davantage de hauteur à sa seconde tentative avant de lâcher prise et de retomber lourdement sur le sol.

Morrÿs essaya à son tour sans parvenir à quitter véritablement le bas de la corde. Azaka, lasse de les regarder lutter, entreprit l'ascension. À mi-hauteur, elle commit l'erreur de regarder vers le bas. Son souffle se bloqua et ses bras se crispèrent autour de la corde.

Kael lui demanda de respirer et de se concentrer sur sa voix. La guide demeura immobile quelques instants, suspendue au-dessus du vide, puis reprit lentement son ascension. Lorsqu'elle atteignit enfin le palier supérieur, elle dut admettre que l'épreuve restait possible. Cela ne signifiait pas qu'elle souhaitait la recommencer.

Obélik ferma les yeux, agrippa la corde et monta sans regarder ni en haut ni en bas. Il parvint au palier suivant dans une obscurité complète, chercha l'épaule d'Azaka à tâtons et ne rencontra que du vide. Morrÿs réussit sa seconde tentative et les rejoignit. Destos échoua à son tour, laissant Kael reprendre son souffle avant un troisième essai.

Le clerc appela sur le rôdeur une brève assistance divine. Cette fois, Kael atteignit le sommet sans retomber. Destos employa ensuite la même prière sur lui-même et se hissa à son tour jusqu'au passage. Lorsque la lumière de son symbole sacré éclaira enfin la cavité, tout le groupe se trouvait réuni au deuxième palier.

\begin{narrateur}
Les araignées avaient été vaincues. La corde, elle, avait opposé une résistance autrement plus sérieuse.
\end{narrateur}

Au-dessus d'eux, Firefinger continuait de brûler dans la nuit. Le masque d'Azaka les attendait encore quelque part dans les hauteurs.
